\subsection{Относительность движения}
%%%%%%%%%%%%%%%%%%%%%%%%%%%%%%%%%%%%%%%%%%%%%%%%%%%%
%%%%%%%%%%%%%%%%%%%%%%%%%%%%%%%%%%%%%%%%%%%%%%%%%%%%
\z{Колонна автомашин длинной $L = 2\, \text{км}$ движется со скоростью $v_1  =36 \, \frac{\text{км}}{\text{ч}}$. Из начала колонны выезжает мотоциклист, который, достигнув ее конца, возвращается обратно. Скорость мотоциклиста постоянная и равна $v_2  = 54\, \frac{\text{км}}{\text{ч}}$. Сколько времени $t$ будет в пути и какой путь $s$ пройдет мотоциклист, пока она снова нагонит начало колонны?}
%
{$t=8\,\text{мин};~s=7.2\, \text{км}$}

%%%%%%%%%%%%%%%%%%%%%%%%%%%%%%%%%%%%%%%%%%%%%%%%%%%%

 \z[{ \begin{wrapfigure}[6]{r}{0.3\textwidth}%
		\vspace{-10pt}%
		\begin{flushright}%
 			\includegraphics[width=0.3\textwidth]{1.4.2}%
		\end{flushright}%
	\end{wrapfigure}%
	}]%
	{Из пункта $A$, расположенного на берегу реки, необходимо попасть в пункт $B$, двигаясь по прямой $AB$ (см. рисунок). Ширина реки $AC = b = 1\, \text{км}$, $BC = a = 2\, \text{км}$, скорость лодки относительно воды $v = 5\, \frac{\text{км}}{\text{ч}}$, а скорость течения реки $u = 2\, \frac{\text{км}}{\text{ч}}$. За какое время $t$ может быть пройдет отрезок $AB$?}
%
{$t=\frac{ua+\sqrt{v^2(a^2+b^2)-u^2b^2}}{v^2-u^2}\approx 43\,\text{мин}$}
	
%%%%%%%%%%%%%%%%%%%%%%%%%%%%%%%%%%%%%%%%%%%%%%%%%%%%
\z{Идет отвесный дождь. Скорость капель $u$. По асфальту со скоростью $v$ скользит мяч. Во сколько раз за один и тот же промежуток времени на него попадает больше капель, чем на такой же, но неподвижный мяч? Изменится ли ответ, если мяч не круглый.}
%
{в $\sqrt{1+\frac{v^2}{u^2}}$, нет}

%%%%%%%%%%%%%%%%%%%%%%%%%%%%%%%%%%%%%%%%%%%%%%%%%%%%	
\z{Минимальное время, которое необходимо, чтобы переплыть в лодке реку, равно $t_0$. Ширина русла реки равна $H$. Скорость течения реки постоянна в любом месте русла и в $\beta$ раз больше скорости лодки ($\beta > 1$), плывущей в стоячей воде.%
   		\lpu Найдите скорость лодки в стоячей воде.\lpuend%
		\lpu На какое расстояние снесёт лодку за минимальное время переправы?\lpuend%
		\lpu Определите наименьшее расстояние, на которое может снести лодку за время переправы.\lpuend%
		\lpu Найдите время переправы лодки в том случае, когда её сносит на минимальное расстояние.\lpuend%
		}
%
{\puan $v=\frac{H}{t_0}$;~\puan $L=\beta H$;~\puan $L_{min}=H\sqrt{\beta^2-1}$;~\puan $t=t_0 \frac{\beta}{\sqrt{\beta^2-1}}$}

%%%%%%%%%%%%%%%%%%%%%%%%%%%%%%%%%%%%%%%%%%%%%%%%%%%%
 \z[{ \begin{wrapfigure}[6]{r}{0.3\textwidth}%
		\vspace{-10pt}%
		\begin{flushright}%
 			\includegraphics[width=0.3\textwidth]{1.4.5}%
		\end{flushright}%
	\end{wrapfigure}%
	}]%
	{По гладкой горизонтальной поверхности стола скользит брусок и ударяет своей гладкой вертикальной гранью $AB$ по шарику, скользящему по столу навстречу бруску (на рисунке показан вид сверху). Скорость бруска составляет угол $\alpha = 60\degree$ с гранью $AB$. После абсолютно упругого удара шарик отскочил со скоростью $v$ под углом $\beta = 45\degree$ к направлению движения бруска. Масса шарика намного меньше массы бруска.%
   		\pu Найдите скорость шарика перед ударом. %
		\pu Найдите скорость бруска.%
	}
%
{\puan $v_0=\frac{\sqrt{3}-1}{\sqrt{2}}v$;~\puan $u=\frac{\sqrt{3}-1}{\sqrt{6}}v$}

